\IfFileExists{stacks-project.cls}{%
\documentclass{stacks-project}
}{%
\documentclass{amsart}
}

% For dealing with references we use the comment environment
\usepackage{verbatim}
\newenvironment{reference}{\comment}{\endcomment}
%\newenvironment{reference}{}{}
\newenvironment{slogan}{\comment}{\endcomment}
\newenvironment{history}{\comment}{\endcomment}

% For commutative diagrams we use Xy-pic
\usepackage[all]{xy}

% We use 2cell for 2-commutative diagrams.
\xyoption{2cell}
\UseAllTwocells

% We use multicol for the list of chapters between chapters
\usepackage{multicol}

% This is generally recommended for better output
\usepackage{lmodern}
\usepackage[T1]{fontenc}

% For cross-file-references
\usepackage{xr-hyper}

% Package for hypertext links:
\usepackage{hyperref}

% For any local file, say "hello.tex" you want to link to please
% use \externaldocument[hello-]{hello}
\externaldocument[introduction-]{introduction}
\externaldocument[conventions-]{conventions}
\externaldocument[sets-]{sets}
\externaldocument[categories-]{categories}
\externaldocument[topology-]{topology}
\externaldocument[sheaves-]{sheaves}
\externaldocument[sites-]{sites}
\externaldocument[stacks-]{stacks}
\externaldocument[fields-]{fields}
\externaldocument[algebra-]{algebra}
\externaldocument[brauer-]{brauer}
\externaldocument[homology-]{homology}
\externaldocument[derived-]{derived}
\externaldocument[simplicial-]{simplicial}
\externaldocument[more-algebra-]{more-algebra}
\externaldocument[smoothing-]{smoothing}
\externaldocument[modules-]{modules}
\externaldocument[sites-modules-]{sites-modules}
\externaldocument[injectives-]{injectives}
\externaldocument[cohomology-]{cohomology}
\externaldocument[sites-cohomology-]{sites-cohomology}
\externaldocument[dga-]{dga}
\externaldocument[dpa-]{dpa}
\externaldocument[sdga-]{sdga}
\externaldocument[hypercovering-]{hypercovering}
\externaldocument[schemes-]{schemes}
\externaldocument[constructions-]{constructions}
\externaldocument[properties-]{properties}
\externaldocument[morphisms-]{morphisms}
\externaldocument[coherent-]{coherent}
\externaldocument[divisors-]{divisors}
\externaldocument[limits-]{limits}
\externaldocument[varieties-]{varieties}
\externaldocument[topologies-]{topologies}
\externaldocument[descent-]{descent}
\externaldocument[perfect-]{perfect}
\externaldocument[more-morphisms-]{more-morphisms}
\externaldocument[flat-]{flat}
\externaldocument[groupoids-]{groupoids}
\externaldocument[more-groupoids-]{more-groupoids}
\externaldocument[etale-]{etale}
\externaldocument[chow-]{chow}
\externaldocument[intersection-]{intersection}
\externaldocument[pic-]{pic}
\externaldocument[weil-]{weil}
\externaldocument[adequate-]{adequate}
\externaldocument[dualizing-]{dualizing}
\externaldocument[duality-]{duality}
\externaldocument[discriminant-]{discriminant}
\externaldocument[derham-]{derham}
\externaldocument[local-cohomology-]{local-cohomology}
\externaldocument[algebraization-]{algebraization}
\externaldocument[curves-]{curves}
\externaldocument[resolve-]{resolve}
\externaldocument[models-]{models}
\externaldocument[functors-]{functors}
\externaldocument[equiv-]{equiv}
\externaldocument[pione-]{pione}
\externaldocument[etale-cohomology-]{etale-cohomology}
\externaldocument[proetale-]{proetale}
\externaldocument[relative-cycles-]{relative-cycles}
\externaldocument[more-etale-]{more-etale}
\externaldocument[trace-]{trace}
\externaldocument[crystalline-]{crystalline}
\externaldocument[spaces-]{spaces}
\externaldocument[spaces-properties-]{spaces-properties}
\externaldocument[spaces-morphisms-]{spaces-morphisms}
\externaldocument[decent-spaces-]{decent-spaces}
\externaldocument[spaces-cohomology-]{spaces-cohomology}
\externaldocument[spaces-limits-]{spaces-limits}
\externaldocument[spaces-divisors-]{spaces-divisors}
\externaldocument[spaces-over-fields-]{spaces-over-fields}
\externaldocument[spaces-topologies-]{spaces-topologies}
\externaldocument[spaces-descent-]{spaces-descent}
\externaldocument[spaces-perfect-]{spaces-perfect}
\externaldocument[spaces-more-morphisms-]{spaces-more-morphisms}
\externaldocument[spaces-flat-]{spaces-flat}
\externaldocument[spaces-groupoids-]{spaces-groupoids}
\externaldocument[spaces-more-groupoids-]{spaces-more-groupoids}
\externaldocument[bootstrap-]{bootstrap}
\externaldocument[spaces-pushouts-]{spaces-pushouts}
\externaldocument[spaces-chow-]{spaces-chow}
\externaldocument[groupoids-quotients-]{groupoids-quotients}
\externaldocument[spaces-more-cohomology-]{spaces-more-cohomology}
\externaldocument[spaces-simplicial-]{spaces-simplicial}
\externaldocument[spaces-duality-]{spaces-duality}
\externaldocument[formal-spaces-]{formal-spaces}
\externaldocument[restricted-]{restricted}
\externaldocument[spaces-resolve-]{spaces-resolve}
\externaldocument[formal-defos-]{formal-defos}
\externaldocument[defos-]{defos}
\externaldocument[cotangent-]{cotangent}
\externaldocument[examples-defos-]{examples-defos}
\externaldocument[algebraic-]{algebraic}
\externaldocument[examples-stacks-]{examples-stacks}
\externaldocument[stacks-sheaves-]{stacks-sheaves}
\externaldocument[criteria-]{criteria}
\externaldocument[artin-]{artin}
\externaldocument[quot-]{quot}
\externaldocument[stacks-properties-]{stacks-properties}
\externaldocument[stacks-morphisms-]{stacks-morphisms}
\externaldocument[stacks-limits-]{stacks-limits}
\externaldocument[stacks-cohomology-]{stacks-cohomology}
\externaldocument[stacks-perfect-]{stacks-perfect}
\externaldocument[stacks-introduction-]{stacks-introduction}
\externaldocument[stacks-more-morphisms-]{stacks-more-morphisms}
\externaldocument[stacks-geometry-]{stacks-geometry}
\externaldocument[moduli-]{moduli}
\externaldocument[moduli-curves-]{moduli-curves}
\externaldocument[examples-]{examples}
\externaldocument[exercises-]{exercises}
\externaldocument[guide-]{guide}
\externaldocument[desirables-]{desirables}
\externaldocument[coding-]{coding}
\externaldocument[obsolete-]{obsolete}
\externaldocument[fdl-]{fdl}
\externaldocument[index-]{index}

% Theorem environments.
%
\theoremstyle{plain}
\newtheorem{theorem}[subsection]{Theorem}
\newtheorem{proposition}[subsection]{Proposition}
\newtheorem{lemma}[subsection]{Lemma}

\theoremstyle{definition}
\newtheorem{definition}[subsection]{Definition}
\newtheorem{example}[subsection]{Example}
\newtheorem{exercise}[subsection]{Exercise}
\newtheorem{situation}[subsection]{Situation}

\theoremstyle{remark}
\newtheorem{remark}[subsection]{Remark}
\newtheorem{remarks}[subsection]{Remarks}

\numberwithin{equation}{subsection}

% Macros
%
\def\lim{\mathop{\mathrm{lim}}\nolimits}
\def\colim{\mathop{\mathrm{colim}}\nolimits}
\def\Spec{\mathop{\mathrm{Spec}}}
\def\Hom{\mathop{\mathrm{Hom}}\nolimits}
\def\Ext{\mathop{\mathrm{Ext}}\nolimits}
\def\SheafHom{\mathop{\mathcal{H}\!\mathit{om}}\nolimits}
\def\SheafExt{\mathop{\mathcal{E}\!\mathit{xt}}\nolimits}
\def\Sch{\mathit{Sch}}
\def\Mor{\mathop{\mathrm{Mor}}\nolimits}
\def\Ob{\mathop{\mathrm{Ob}}\nolimits}
\def\Sh{\mathop{\mathit{Sh}}\nolimits}
\def\NL{\mathop{N\!L}\nolimits}
\def\CH{\mathop{\mathrm{CH}}\nolimits}
\def\proetale{{pro\text{-}\acute{e}tale}}
\def\etale{{\acute{e}tale}}
\def\QCoh{\mathit{QCoh}}
\def\Ker{\mathop{\mathrm{Ker}}}
\def\Im{\mathop{\mathrm{Im}}}
\def\Coker{\mathop{\mathrm{Coker}}}
\def\Coim{\mathop{\mathrm{Coim}}}

% Boxtimes
%
\DeclareMathSymbol{\boxtimes}{\mathbin}{AMSa}{"02}

%
% Macros for moduli stacks/spaces
%
\def\QCohstack{\mathcal{QC}\!\mathit{oh}}
\def\Cohstack{\mathcal{C}\!\mathit{oh}}
\def\Spacesstack{\mathcal{S}\!\mathit{paces}}
\def\Quotfunctor{\mathrm{Quot}}
\def\Hilbfunctor{\mathrm{Hilb}}
\def\Curvesstack{\mathcal{C}\!\mathit{urves}}
\def\Polarizedstack{\mathcal{P}\!\mathit{olarized}}
\def\Complexesstack{\mathcal{C}\!\mathit{omplexes}}
% \Pic is the operator that assigns to X its picard group, usage \Pic(X)
% \Picardstack_{X/B} denotes the Picard stack of X over B
% \Picardfunctor_{X/B} denotes the Picard functor of X over B
\def\Pic{\mathop{\mathrm{Pic}}\nolimits}
\def\Picardstack{\mathcal{P}\!\mathit{ic}}
\def\Picardfunctor{\mathrm{Pic}}
\def\Deformationcategory{\mathcal{D}\!\mathit{ef}}
\usepackage{cleveref}
\crefname{originalequation}{original equation}{original equations}
\Crefname{originalequation}{Original equation}{Original equations}
\makeatletter
\newlabel{equation-decompose}{{015T}{source}{}{sourcepointer-1}{}}
\newlabel{equation-decompose@cref}{{[originalequation][0][]015T}{[1][0][]source}{}{}{}}
\makeatother
\crefname{originalequation}{original equation}{original equations}
\Crefname{originalequation}{Original equation}{Original equations}
\makeatletter
\newlabel{equation-definition-filtered-derived-functor}{{05U2}{source}{}{sourcepointer-2}{}}
\newlabel{equation-definition-filtered-derived-functor@cref}{{[originalequation][0][]05U2}{[1][0][]source}{}{}{}}
\makeatother
\crefname{originalreference}{original reference}{original references}
\Crefname{originalreference}{Original reference}{Original references}
\makeatletter
\newlabel{homology-lemma-biregular-ss-converges}{{012W}{source}{}{sourcepointer-3}{}}
\newlabel{homology-lemma-biregular-ss-converges@cref}{{[originalreference][0][]012W}{[1][0][]source}{}{}{}}
\makeatother
\crefname{originalreference}{original reference}{original references}
\Crefname{originalreference}{Original reference}{Original references}
\makeatletter
\newlabel{homology-lemma-spectral-sequence-filtered-complex}{{012M}{source}{}{sourcepointer-4}{}}
\newlabel{homology-lemma-spectral-sequence-filtered-complex@cref}{{[originalreference][0][]012M}{[1][0][]source}{}{}{}}
\makeatother
\crefname{originalreference}{original reference}{original references}
\Crefname{originalreference}{Original reference}{Original references}
\makeatletter
\newlabel{homology-lemma-spectral-sequence-filtered-complex-functorial}{{012O}{source}{}{sourcepointer-5}{}}
\newlabel{homology-lemma-spectral-sequence-filtered-complex-functorial@cref}{{[originalreference][0][]012O}{[1][0][]source}{}{}{}}
\makeatother
\crefname{originalreference}{original reference}{original references}
\Crefname{originalreference}{Original reference}{Original references}
\makeatletter
\newlabel{homology-section-filtered-complex}{{012K}{source}{}{sourcepointer-6}{}}
\newlabel{homology-section-filtered-complex@cref}{{[originalreference][0][]012K}{[1][0][]source}{}{}{}}
\makeatother
\crefname{originallemma}{original lemma}{original lemmas}
\Crefname{originallemma}{Original lemma}{Original lemmas}
\makeatletter
\newlabel{lemma-morphisms-into-filtered-injective-complex}{{05TY}{source}{}{sourcepointer-7}{}}
\newlabel{lemma-morphisms-into-filtered-injective-complex@cref}{{[originallemma][0][]05TY}{[1][0][]source}{}{}{}}
\makeatother
\crefname{originallemma}{original lemma}{original lemmas}
\Crefname{originallemma}{Original lemma}{Original lemmas}
\makeatletter
\newlabel{lemma-right-resolution-by-filtered-injectives}{{05TW}{source}{}{sourcepointer-8}{}}
\newlabel{lemma-right-resolution-by-filtered-injectives@cref}{{[originallemma][0][]05TW}{[1][0][]source}{}{}{}}
\makeatother
\title{The filtered spectral sequence for a left exact functor on filtered complexes}
\author{The Stacks Project Authors}
\date{stacks-project:derived.tex:local-snapshot-20261012; GNU FDL 1.2 or later}
\begin{document}
\maketitle
\noindent\textit{Editorial note.} Counted span is the author's own proof environment, source lines 8722-8773, following the statement at 8700-8720.\par\medskip
\noindent\textit{Editorial note.} No omissions: every equality in the E\_1 chain is displayed with its reason.\par\medskip
\noindent\textit{Allowed standard prerequisites (editorial).}\begin{itemize}
\item Filtered injective resolution. lemma-right-resolution-by-filtered-injectives (carded L26168), cited at proof lines 8725 and 8756. Supplies the bounded below filtered injective resolutions K -> I and L -> J.
\item Spectral sequence of a filtered complex. homology-lemma-spectral-sequence-filtered-complex and its functorial variant, cross-chapter citations at proof lines 8733 and 8771 (tags mechanism). Supplies E\_1\textasciicircum{}\{p,q\} = H\textasciicircum{}\{p+q\}(gr\textasciicircum{}p(T(I))) and the functorial morphism of spectral sequences.
\item Biregular spectral sequences converge. homology-lemma-biregular-ss-converges, cross-chapter citation at proof line 8745 (tags mechanism). Supplies boundedness, convergence to R\textasciicircum{}n T(K) and the finite filtrations.
\item Maps into filtered injective complexes. lemma-morphisms-into-filtered-injective-complex, cited at proof line 8759; complete proof at source lines 8482-8523, queued for separate carding. Supplies the homotopy-commutative square between the two resolutions used for functoriality.
\end{itemize}\medskip
\ifdefined\hypertarget\hypertarget{sourcepointer-1}{}\fi
\noindent\textit{Source reference.} Original equation 015T: sources/stacks-derived/pinned-original-source.tex, source lines 8635--8635. Navigation pointer only: the counted proof cites this statement by \textbackslash{}ref as a standing input; the deduction the counted proof has to make is carried out inside the counted span, and no part of that deduction is delegated to this pointer.\par\medskip
\ifdefined\hypertarget\hypertarget{sourcepointer-2}{}\fi
\noindent\textit{Source reference.} Original equation 05U2: sources/stacks-derived/pinned-original-source.tex, source lines 8659--8659. Navigation pointer only: the counted proof cites this statement by \textbackslash{}ref as a standing input; the deduction the counted proof has to make is carried out inside the counted span, and no part of that deduction is delegated to this pointer.\par\medskip
\ifdefined\hypertarget\hypertarget{sourcepointer-3}{}\fi
\noindent\textit{Source reference.} Original reference 012W: focused-seeds/geometry/S9-der-ss-filtered-derived/ref-tags.txt, source lines 5--5. Navigation pointer only: the counted proof cites this statement by \textbackslash{}ref as a standing input; the deduction the counted proof has to make is carried out inside the counted span, and no part of that deduction is delegated to this pointer.\par\medskip
\ifdefined\hypertarget\hypertarget{sourcepointer-4}{}\fi
\noindent\textit{Source reference.} Original reference 012M: focused-seeds/geometry/S9-der-ss-filtered-derived/ref-tags.txt, source lines 6--6. Navigation pointer only: the counted proof cites this statement by \textbackslash{}ref as a standing input; the deduction the counted proof has to make is carried out inside the counted span, and no part of that deduction is delegated to this pointer.\par\medskip
\ifdefined\hypertarget\hypertarget{sourcepointer-5}{}\fi
\noindent\textit{Source reference.} Original reference 012O: focused-seeds/geometry/S9-der-ss-filtered-derived/ref-tags.txt, source lines 7--7. Navigation pointer only: the counted proof cites this statement by \textbackslash{}ref as a standing input; the deduction the counted proof has to make is carried out inside the counted span, and no part of that deduction is delegated to this pointer.\par\medskip
\ifdefined\hypertarget\hypertarget{sourcepointer-6}{}\fi
\noindent\textit{Source reference.} Original reference 012K: focused-seeds/geometry/S9-der-ss-filtered-derived/ref-tags.txt, source lines 8--8. Navigation pointer only: the counted proof cites this statement by \textbackslash{}ref as a standing input; the deduction the counted proof has to make is carried out inside the counted span, and no part of that deduction is delegated to this pointer.\par\medskip
\ifdefined\hypertarget\hypertarget{sourcepointer-7}{}\fi
\noindent\textit{Source reference.} Original lemma 05TY: sources/stacks-derived/pinned-original-source.tex, source lines 8458--8458. Navigation pointer only: the counted proof cites this statement by \textbackslash{}ref as a standing input; the deduction the counted proof has to make is carried out inside the counted span, and no part of that deduction is delegated to this pointer.\par\medskip
\ifdefined\hypertarget\hypertarget{sourcepointer-8}{}\fi
\noindent\textit{Source reference.} Original lemma 05TW: sources/stacks-derived/pinned-original-source.tex, source lines 8313--8313. Navigation pointer only: the counted proof cites this statement by \textbackslash{}ref as a standing input; the deduction the counted proof has to make is carried out inside the counted span, and no part of that deduction is delegated to this pointer.\par\medskip
\begin{lemma}
\label{lemma-ss-filtered-derived}
Let $\mathcal{A}, \mathcal{B}$ be abelian categories. Let
$T : \mathcal{A} \to \mathcal{B}$ be a left exact functor.
Assume $\mathcal{A}$ has enough injectives.
Let $(K^\bullet, F)$ be an object of
$\text{Comp}^{+}(\text{Fil}^f(\mathcal{A}))$.
There exists a spectral sequence $(E_r, d_r)_{r\geq 0}$
consisting of bigraded objects $E_r$ of $\mathcal{B}$
and $d_r$ of bidegree $(r, - r + 1)$ and with
$$
E_1^{p, q} = R^{p + q}T(\text{gr}^p(K^\bullet))
$$
Moreover, this spectral sequence is bounded, converges
to $R^*T(K^\bullet)$, and induces a finite
filtration on each $R^nT(K^\bullet)$. The construction
of this spectral sequence is functorial in the object
$K^\bullet$ of $\text{Comp}^{+}(\text{Fil}^f(\mathcal{A}))$
and the terms $(E_r, d_r)$ for $r \geq 1$ do not depend
on any choices.
\end{lemma}

\begin{proof}
Choose a filtered quasi-isomorphism $K^\bullet \to I^\bullet$
with $I^\bullet$ a bounded below complex of filtered injective objects, see
Lemma \ref{lemma-right-resolution-by-filtered-injectives}.
Consider the complex $RT(K^\bullet) = T_{ext}(I^\bullet)$, see
(\ref{equation-definition-filtered-derived-functor}).
Thus we can consider the spectral sequence
$(E_r, d_r)_{r \geq 0}$ associated to
this as a filtered complex in $\mathcal{B}$, see
Homology, Section \ref{homology-section-filtered-complex}.
By
Homology, Lemma \ref{homology-lemma-spectral-sequence-filtered-complex}
we have $E_1^{p, q} = H^{p + q}(\text{gr}^p(T(I^\bullet)))$.
By Equation (\ref{equation-decompose}) we have
$E_1^{p, q} = H^{p + q}(T(\text{gr}^p(I^\bullet)))$, and
by definition of a filtered injective resolution the
map $\text{gr}^p(K^\bullet) \to \text{gr}^p(I^\bullet)$
is an injective resolution. Hence
$E_1^{p, q} = R^{p + q}T(\text{gr}^p(K^\bullet))$.

\medskip\noindent
On the other hand, each $I^n$ has a finite filtration and hence
each $T(I^n)$ has a finite filtration. Thus we may apply
Homology, Lemma \ref{homology-lemma-biregular-ss-converges}
to conclude that the spectral sequence is bounded, converges to
$H^n(T(I^\bullet)) = R^nT(K^\bullet)$
moreover inducing finite filtrations on each of the terms.

\medskip\noindent
Suppose that $K^\bullet \to L^\bullet$ is a morphism of
$\text{Comp}^{+}(\text{Fil}^f(\mathcal{A}))$.
Choose a filtered quasi-isomorphism $L^\bullet \to J^\bullet$
with $J^\bullet$ a bounded below complex of filtered injective
objects, see
Lemma \ref{lemma-right-resolution-by-filtered-injectives}.
By our results above,
for example
Lemma \ref{lemma-morphisms-into-filtered-injective-complex},
there exists a diagram
$$
\xymatrix{
K^\bullet \ar[r] \ar[d] & L^\bullet \ar[d] \\
I^\bullet \ar[r] & J^\bullet
}
$$
which commutes up to homotopy. Hence we get a morphism of filtered
complexes $T(I^\bullet) \to T(J^\bullet)$ which gives rise to the
morphism of spectral sequences, see
Homology,
Lemma \ref{homology-lemma-spectral-sequence-filtered-complex-functorial}.
The last statement follows from this.
\end{proof}

\end{document}
